% !TEX program = xelatex
\documentclass{英语试卷}
\begin{document}
\examtitle{惠州市 2025 届高三第二次调研考试试题}{英语}
\par 惠州市 2025 届高三第二次调研考试试题\par
\begin{notice}
\par 1. 答卷前，考生务必用黑色字迹的钢笔或签字笔将自己的姓名、考生号、试室号和座位号填写在答题卡上。用 2B 铅笔将试卷类型（A）填涂在答题卡相应位置上，并在答题卡相应位置上填涂考生号。因笔试不考听力，试卷从第二部分开始，试题序号从“21"开始。\par
\par 2. 作答选择题时，选出每小题答案后，用 2B 铅笔把答题卡对应题目选项的答案信息点除黑；如需改动，用橡皮擦干净后，再选涂其他答案。答案不能答在试卷上，\par
\par 3. 非选择题必须用黑色字迹的钢笔或签字笔作答，答案必须写在答题卡各题目指定区域内相应位置上，如需改动，先划掉原来的答案，然后再写上新答案；不准使用铅笔和涂改液。不按以上要求作答无效。\par
\par 4. 考生必须保持答题卡的整洁。考试结束后，将试卷和答题卡一并交回，\par
\end{notice}
\parttitle{第二部分 阅读理解（共两节，满分50分）}
\sectitle{第一节（共15小题；每小题2.5分，满分37.5分）}
\par 阅读下列短文，从每题所给的 A、B、C 和 D 四个选项中，选出最佳答案。\par
\passagelabel{A}
\par Fall Festival: Fall Reimagined at the Detroit Zoo\par
\par This October, FALL for the Zoo in a whole new way! We're bringing you a variety of events that promise something for everyone—from family-friendly fun to 21-and-older happenings.\par
\par Spooky Spirits\par
\par Experience exciting entertainment at Spooky Spirits, our tasting event for those aged 21 and over. As a bonus, you'll get a special first look at our newest feature—Jack O' Lantern Journey! With more than 80 drinks to choose from and delicious food options, fall has never tasted so good.\par
\par Jack O' Lantern Journey\par
\par This journey features illuminated(发光的) pumpkin displays, and photo opportunities where the magic of the season comes to life. Carve out new memories as you experience the BRAND-NEW Jack O' Lantern Journey! This must-see experience is perfect for a date night, family outing or night out with your friends.\par
\par Zoo Boo\par
\par Explore the Zoo, experience festive pumpkin displays, and enjoy strolling entertainment. You can capture memories with themed photo opportunities while visiting each of the treat station spots throughout the Zoo. The Zoo will be open until sunset, allowing you to interact with your favorite animals while celebrating the season. This event is perfect for families with young children!\par
\par Blake's Market\par
\par Experience the true essence of autumn at Blake's Market—from pumpkins to irresistible pies. At the Detroit Zoo's Blake's Market, you'll find tasty treats and the perfect pumpkin for all your carving needs. You are expected to enjoy delicious seasonal food and baked goods for sale.\par
\par\vspace{0.4em}
\begin{center}\begin{tabular}{{|l|l|l|}}
\hline
Event & Time & Price \\
Spooky Spirits & Oct. 4th & \$ 40 \\
Jack O' Lantern Journey & Oct. \{5\}th - \{6\}th & \$ 19 \\
Zoo Boo & Oct. \{11\}th - \{18\}th & \$ 19 \\
Blake's Market & Oct. \{20\}th - \{26\}th & Free \\
\hline
\end{tabular}\end{center}
\vspace{0.4em}
\question{21. Which event is specially intended for grown-ups?}
\begin{choices}(2)
  \item Spooky Spirits.
  \item Jack O' Lantern Journey.
  \item Zoo Boo.
  \item Blake's Market.
\end{choices}
\question{22. What can visitors do on Oct.17?}
\begin{choices}(2)
  \item Make pumpkin pies.
  \item Visit their favorite animals.
  \item Enjoy various drinks.
  \item Take pictures of wildlife.
\end{choices}
\question{23. What is the main objective of the Fall Festival at the Detroit Zoo?}
\begin{choices}(2)
  \item To provide a platform for local craftsmen.
  \item To introduce the zoo's newest animal exhibits.
  \item To celebrate the autumn season and attract visitors.
  \item To raise funds for the conservation efforts of the zoo.
\end{choices}
\passagelabel{B}
\par On a windy day, 40 middle-school students stand by the Harlem River in New York City. Their teacher, Mr. White, pulls a metal cage out of the river. Inside are clusters of oysters(牡蛎). The last time they were here, the oysters were babies.\par
\par The students' job is to measure how much the oysters have grown and test the river water for pollutants. They will report their findings to researchers at the Billion Oyster Project (BOP). Why work so hard? To make the city's waterways healthy again.\par
\par Oysters are important for the underwater community they share with plants, fish, and other life. So they are also called a keystone species. If a keystone species disappears, other plants and animals may die out. That's exactly what happened 100 years ago in New York Harbor. Before then, lots of oysters lived here. The oysters were delicious—maybe too delicious. By the early 1900s, people were eating them faster than they could grow.\par
\par The oysters were in trouble for another reason. Pollution was pouring into the waters of the growing city, and few types of animals could live in it. Since the 1970s, new laws have helped reduce toxic(有毒的) waste. Some fish started to swim through again. But oysters were still missing—until recently. The BOP began in 2014 to help bring oysters back to New York Harbor. So far, students and volunteers have placed about 25,000,000 oysters into the harbor.\par
\par "The oysters definitely look much bigger now," says a student named Kelly. Taking turns, the kids measure all the oysters, then compare notes. The biggest oyster is over 2 inches long, much bigger than a healthy size for its age.\par
\par As more oysters grow, the water should become clearer. Also, other animals should move in. Finally, the students put the oysters back in the cage. Mr. White lowers it into the river. When the oysters are big enough, they will be moved to the middle of the harbor. For today, everyone agrees: the future is looking bright for New York Harbor.\par
\question{24. Why are oysters called a keystone species?}
\begin{choices}(2)
  \item They grow quickly.
  \item They are delicious.
  \item They support underwater life.
  \item They are able to resist pollution.
\end{choices}
\question{25. What was the problem with oysters over a century ago?}
\begin{choices}(2)
  \item A period of disappearance.
  \item Over-consumption by people.
  \item Lack of clean and safe habitat.
  \item Competition with other species.
\end{choices}
\question{26. What's the goal of the Billion Oyster Project?}
\begin{choices}(2)
  \item To grow more oysters.
  \item To collect useful data.
  \item To protect water plants.
  \item To restore waterways.
\end{choices}
\question{27. What may be the future outlook for New York Harbor?}
\begin{choices}(2)
  \item Promising.
  \item Uncertain.
  \item Disappointing.
  \item Unchanged.
\end{choices}
\passagelabel{C}
\par Traditional micro-robots often face limitations due to their power sources, but recently developed insect-sized robot powered by tiny explosions has demonstrated the ability to climb, jump, and carry loads many times its own weight.\par
\par The robot, mainly developed by materials engineer Robert Shepherd at Cornell University in New York, along with his PhD student Cameron Aubin, is powered by tiny actuators(致动器). "The actuator looks like a column. It’s an empty cylinder with a stretchy rubber on the top," said Aubin.\par
\par The researchers employed four actuators to drive the robot's feet. To make the robot jump or climb, a stream of methane (甲烷) and oxygen was fed into each foot and sparked (触发) with electricity from a battery. The resulting reaction between the gases released energy as a small explosion, causing the rubber to change its form.\par
\par Shepherd noted that the rapid explosions were so tiny that they didn't damage the rubber. But they provided considerable propulsion: the robot could jump to a height of 56 centimeters and carry a load 22 times its own weight.\par
\par "This team has brought chemically powered actuation to impressive ranges for robotics while also demonstrating impressive capabilities for insect-size machines," said Ryan Truby, a materials scientist at Northwestern University in Illinois.\par
\par Robots that are small and light but strong and able to cover large distances could one day be used in environmental monitoring or search-and-rescue applications. And the robot built by Aubin and Shepherd, described in a paper published in Science, would need further development for use in the field. "At the moment, it needs to be connected to a panel that contains the gas supply and a battery, a limitation that presents a challenge. Coming up with ways to burn fuels on a wireless robot is not easy," said Truby. "This is the next big challenge to address."\par
\question{28. Why does the recently developed insect- sized robot stand out?}
\begin{choices}(2)
  \item It is extremely tiny and light.
  \item It is equipped with a column.
  \item It is extremely swift and quick.
  \item It is driven by explosive force.
\end{choices}
\question{29. What happens when the methane and oxygen is sparked?}
\begin{choices}(2)
  \item The reaction creates a mixture of gas to drive the robot.
  \item The reaction generates heat that is changed into movement.
  \item The reaction transforms a stretchy rubber to drive the robot.
  \item The reaction provides electricity that is changed into movement.
\end{choices}
\question{30. What does "propulsion" in paragraph 4 probably mean?}
\begin{choices}(2)
  \item Ability to move.
  \item Possibility to develop.
  \item Force to push things.
  \item Electricity to spark gases.
\end{choices}
\question{31. How can Aubin and Shepherd make advancement to the robot they built?}
\begin{choices}(2)
  \item They can use more actuators to provide larger power.
  \item They can design programs to control the robot efficiently.
  \item They can work on methods to enable energy production wirelessly.
  \item They can connect the robot to a panel containing gas supply and a battery.
\end{choices}
\passagelabel{D}
\par Researchers at the Potsdam Institute for Climate Impact Research analyzed historical food prices across various food categories and countries globally, looking at the impact of climatic changes on food inflation(通货膨胀) and the potential implications of future climate change. They found that annual food inflation will rise by 3.2 percent per year and overall inflation by 1.18 percent per year by 2035. And all nations will be impacted, while hot regions and summers will see the most significant consequences.\par
\par "Looking at over 27,000 observations of historical data, we found that increases in temperatures can increase food prices, particularly in hot regions and seasons," said Max Kotz, a co-author of the study, "Under future climate conditions, these impacts could become large, threatening the price stability commands of central banks, such as the ECB (European Central Bank), which aims to keep inflation below 2 percent."\par
\par The report highlighted how rising or unstable prices end anger human welfare and economic and political stability, pointing to United Nations data showing that the cost-of-living crisis in 2021/22 resulted in 71 million more people falling into poverty globally. The study found that increases in temperatures, particularly in hot regions and during summers, lead to heightened inflation levels, which are more pronounced in lower latitudes(纬度) like Africa and South America. Regions at higher latitudes experience seasonal peaks, mostly in the summer months.\par
\par Data showed a 1°C rise in average monthly temperatures affects prices for nearly a year, similarly to extra rainfall. However, price impacts resulting from extreme drought are only temporary. Researchers examined the extremely hot European summer of 2022, when heat and drought had far-reaching effects on agriculture and the economy. "We estimate that the 2022 summer heat extreme increased food inflation in Europe by about 0.6 percent," said Kotz.\par
\par Researchers claimed that future warming projected for 2035 would multiply the impacts of such extremes by up to 50 percent. These effects are particularly relevant for currency unions with a 2 percent inflation target, such as the Eurozone, and will continue to increase with future global warming.\par
\question{32. What is the first paragraph mainly about?}
\begin{choices}(2)
  \item Natural consequences of climate change.
  \item The link between food prices and climate.
  \item Food prices in hot regions during summers.
  \item The impact of food inflation on food prices.
\end{choices}
\question{33. What might the rising prices lead to according to the researchers' report?}
\begin{choices}(2)
  \item Massive poverty.
  \item Human welfare.
  \item Environmental crisis.
  \item Economic stability.
\end{choices}
\question{34. Why is the European summer of 2022 mentioned in paragraph 4?}
\begin{choices}(2)
  \item To give an example of the research findings.
  \item To provide reasons for the researchers' doubt.
  \item To show practical significance of the research.
  \item To state researchers' methods for data analysis.
\end{choices}
\question{35. What is the researchers' attitude toward food inflation?}
\begin{choices}(2)
  \item Unclear.
  \item Dismissive.
  \item Doubtful.
  \item Concerned.
\end{choices}
\sectitle{第二节（共5小题；每小题2.5分，满分12.5分）}
\par 根据短文内容，从短文后的选项中选出能填入空白处的最佳选项。选项中有两项为多余选项。\par
\par Launch Your Spiritual Journey\par
\par A spiritual journey is a personal adventure where you divide deep into understanding and connecting with your inner self. \blanknum{36} To achieve a profound sense of your spiritual essence. Here's a simple guide on how to start your spiritual journey.\par
\par Set your intentions\par
\par Like preparing for a marathon, the first step is to establish your goals. \blanknum{37} Are you looking for inner peace or a deeper understanding of your life purpose? As you navigate your spiritual journey, you can reflect on these intentions, reminding yourself of why you started and what you're striving for.\par
\par Read and learn\par
\par Consider books about spirituality and personal growth as your guidance. They are packed with valuable insights and practical advice. Keep in mind, it's not about gathering much information. \blanknum{38} So take moments to think about these ideas and see how they fit with your experiences.\par
\par Find a mentor\par
\par A mentor can serve a significant role in your spiritual exploration. And your mentor could be anyone----a spiritual teacher, a trusted friend, or even a family member, who has a profound understanding of the spiritual world. They've journeyed through the ups and downs, faced challenges, and emerged enlightened. \blanknum{39} It can brighten your path and make the journey feel less overwhelming.\par
\par Maintain a spiritual diary\par
\par This diary is a tangible reflection of your inner explorations, and plays a vital role in your spiritual progression. Writing regularly can foster mindfulness and enhance self-awareness. \blanknum{40} Finally, this diary can allow you to revisit those past entries, recognize your growth, and draw inspiration for the future.\par
\par Now that you've uncovered the steps to launch a spiritual journey, why not get a taste of this transformative experience?\par
\begin{choices}(1)
  \item But the truth?
  \item And the ultimate goal?
  \item Their wisdom, therefore, can be a lighthouse for you.
  \item It's a voyage into the deepest corners of your memories.
  \item This means identifying what you hope to achieve or discover.
  \item It's about finding the wisdom that truly speaks to your journey.
  \item It can also help you express and strengthen your understandings.
\end{choices}
\parttitle{第三部分 语言知识运用（共两节，满分 30 分）}
\sectitle{第一节（共15小题；每小题1分，满分15分）}
\par 阅读下面短文，从短文后各题所给的A、B、C和D四个选项中，选出可以填入空白处的最佳选项。\par
\par I wasn't an outgoing kid and just didn't dare to say much. I neither got in trouble nor got \blanknum{41}. My classmates would call me a nerd (书呆子). Somehow, Mr. Dalton's plane geometry class really clicked for me.\par
\par Toward the end of the year, Mr. Dalton wanted to check our \blanknum{42} of geometry and gave us a 25-question true/false test. The test presented each \blanknum{43} with a conclusion shown. If we could \blanknum{44} the result using our theorems(定理), the answer was true; otherwise, it was false.\par
\par I breezed through 24 problems, but number 17 was a \blanknum{45}! I spent 30 minutes on that proof, \blanknum{46} my page with theorems and calculations, until the proof was there at last. \blanknum{47}, I marked the question true and handed in the test.\par
\par \blanknum{48}, I received a 96, and number 17 was marked wrong! When Mr. Dalton asked if there were any questions, my hand \blanknum{49}. "The correct answer to number 17 is true!"\par
\par Instead of \blanknum{50} that this was a nationally standardized test and what I could possibly know about the answer, Mr. Dalton gave me an \blanknum{51} smile and asked me to present my proof to the class. After my demonstration, everyone agreed, with \blanknum{52} no less, that the correct answer was true.\par
\par Ever since that geometry class, I've developed the confidence to \blanknum{53} my own ideas. I am grateful to Mr. Dalton for allowing me to challenge his \blanknum{54}. The experience of that day was like a \blanknum{55} from Mr. Dalton and it has served me well in my life.\par
\cz{41}{noticed}{caught}{heard}{ignored}
\cz{42}{intelligence}{method}{awareness}{knowledge}
\cz{43}{answer}{problem}{choice}{theory}
\cz{44}{correct}{prove}{change}{acquire}
\cz{45}{sign}{mistake}{challenge}{chance}
\cz{46}{decorating}{wrapping}{combining}{covering}
\cz{47}{Confused}{Satisfied}{Frustrated}{Astonished}
\cz{48}{Doubtfully}{Fortunately}{Surprisingly}{Excitedly}
\cz{49}{drew back}{reached out}{shot up}{hang down}
\cz{50}{complaining}{arguing}{discussing}{persuading}
\cz{51}{encouraging}{embarrassing}{uncomfortable}{innocent}
\cz{52}{laughter}{pride}{relief}{applause}
\cz{53}{voice}{reflect}{find}{change}
\cz{54}{responsibility}{reputation}{patience}{authority}
\cz{55}{gift}{talent}{goal}{dream}
\sectitle{第二节（共10小题；每小题1.5分，满分15分）}
\par 阅读下面短文，在空白处填入 1 个适当的单词或括号内单词的正确形式。\par
\par Inspired by Marco Polo, Vienna Cammarota set off from the legendary explorer's hometown of Venice to retrace the route on foot. Beginning in Venice on April 26, 2022, \blanknum{} well-known hiker in her 70s arrived in China on August 29, 2024, covering more than 20,000 kilometers.\par
\par Cammarota has long been an adventurous and culturally driven traveler. When \blanknum{57} (ask) why she challenged herself \blanknum{58} such an old age, she answered, "Each new journey brings me fresh experiences and \blanknum{59} (memory)." During this journey, many kindhearted people in countries involved in the Belt and Road Initiative \blanknum{60} (help) her, making her long journey more fulfilling and \blanknum{61} (enjoy).\par
\par The moment she set foot in China, Cammarota was \blanknum{62} (warm) greeted by people who presented her with a wreath, symbolizing friendship and blessings. \blanknum{63} (start) from the ancient city of Kashgar in the autonomous region, Cammarota will embark on a 15-month cultural journey across China and explore over 30 cities in 17 provinces, through \blanknum{64} she will showcase China's true image, rich culture and social harmony to the world. It \blanknum{65} (believe) that her experiences will not only attract more international tourists to China, but also promote cultural dissemination(传播) and economic development in cities involved in the BRI.\par
\parttitle{第四部分 写作（共两节，满分 40 分）}
\sectitle{第一节（满分15分）}
\par 假定你是李华,你校举办了一场以 "Green Life, Healthy Plogging" 为主题的志愿者活动。你的外国笔友 Mike 对此很感兴趣，向你了解相关情况，请你给他回信，内容包括：\par
\par (1) 活动目的;\par
\par (2) 活动内容;\par
\par (3) 你的感想。\par
\par 参考词汇: plogging 捡跑(在慢跑过程中捡拾垃圾的运动) 注意:\par
\par (1) 写作词数应为 80 左右;\par
\par (2) 请按如下格式在答题卡的相应位置作答。\par
\par\vspace{0.4em}
\begin{center}\begin{tabular}{{|l|}}
\hline
Dear Mike, \\
Yours, Li Hua \\
\hline
\end{tabular}\end{center}
\vspace{0.4em}
\sectitle{第二节（满分25分）}
\par 阅读下面材料，根据其内容和所给段落开头语续写两段，使之构成一篇完整的短文。\par
\par The Kids' Table\par
\par James was exhausted. Three families, including four kids under age six, were crowded into his family's house for Thanksgiving.\par
\par "James! James!" called a little voice. Oh no! Three-year-old Cara had found James's hiding spot. James was just trying to catch his breath. All day, he'd been grabbed and bumped by little cousins who thought he was a climbing gym. Cara now leaped onto James' back and held tightly to his neck. James emerged in the kitchen where the adults were getting dinner ready.\par
\par "OK, time to let go. You're choking me," James gasped as he set Cara on the floor. Suddenly, Mario, Dana and Kyle stormed into the kitchen, shouting. "I want a piggyback, too!" "No, it's my turn!"\par
\par Fortunately, Mom saved him. "No, it's not time to play with James! It's almost time for dinner. Go wash your hands."\par
\par "Awww," whined(哭嚷) the kids. Mom shooed(赶走) them out and then stroked James' hair, "No rest for the main attraction, eh?" James ducked(躲避). He hated being treated like a little kid. He was almost as tall as his mom, after all. Mom handed him a stack of plastic plates. "Here, take these to the kids’table."\par
\par The kids' table. He'd forgotten that he would be eating with the babies, separate from the adults. He hoped the kids would be too hungry to bother him. As soon as he set down the last plate, the kids charged back into the room. James sighed when Mom patted his shoulder. "Maybe you'd like to sit with the adults this year," she said. "I saved you a place next to Dad. I'll sit at the kids' table."\par
\par The adults' table! No screaming, no juice spilled on his lap, no dirty fingers stealing his bread. James sat next to his father, too happy to speak. The conversation around him centered on electric bills, how to repair a roof, and so on.\par
\par 注意:\par
\par 1. 续写词数应为 150 左右;\par
\par 2. 请按如下格式在答题卡的相应位置作答，\par
\writingpara{Paragraph 1: Yet, James just couldn't help but glance at the kids.}
\writinglines{8}
\writingpara{Paragraph 2: Seeing Mom struggling to handle the chaos, James went to the kids table.}
\writinglines{8}
\end{document}
